% !PRINTED none
% Section-opening page: no running page number visible in the scan; part title, section
% heading, and body text follow.

\begin{center}
\gesperrt{Zwölfter Abschnitt.}

\vspace{1em}
\gesperrt{Theorie der Einheitswurzeln.}

\vspace{0.5em}
\noindent\rule{4cm}{0.4pt}

\vspace{1em}
\S.\,131.

\vspace{0.5em}
\gesperrt{Die Einheitswurzeln.}
\end{center}

\vspace{1em}
Unter einer Einheitswurzel versteht man allgemein eine
reelle oder imaginäre Zahl von der Beschaffenheit, dass irgend
eine ihrer Potenzen mit einem ganzzahligen Exponenten gleich
1 ist. Diese Einheitswurzeln haben wir schon im \S.\,33 kennen
gelernt und durch trigonometrische und Exponentialfunctionen
dargestellt. Wir haben ferner specielle Fälle, z.\,B. im \S.\,35 die
dritten Einheitswurzeln, benutzt, die dort ohne trigonometrische
Functionen ausgedrückt wurden.

In allen tiefer gehenden Untersuchungen über algebraische
Gleichungen ist nun eine genauere Kenntniss der Einheitswurzeln
und ihrer Eigenschaften unerlässlich. Wir werden uns daher in
diesem Abschnitt mit dem elementaren Theil der algebraischen
Theorie der Einheitswurzeln eingehender beschäftigen, ohne von
ihrer Darstellung durch trigonometrische Functionen Gebrauch
zu machen. Wegen der geometrischen Anwendung auf die Construction
der regulären Vielecke, auf die wir schon im \S.\,33
hingewiesen haben, wird die Theorie der Einheitswurzeln auch
die \gesperrt{Kreistheilungstheorie} genannt. Sie ist im Wesentlichen
eine Schöpfung von \gesperrt{Gauss} $^1$).

\noindent\rule{4cm}{0.4pt}

$^1$) \gesperrt{Gauss}, Disq. ar. Sectio VII.
